\documentclass[a4paper,12pt]{article}

% --- Use XeLaTeX font support ---

\usepackage{fontspec}
\setmainfont{TeX Gyre Termes} % nice Times-like font, installed by texlive-core
\setsansfont{TeX Gyre Heros} % optional sans-serif
\setmonofont{Fira Code}       % optional monospaced font
% --- Math packages ---
\usepackage{amsmath, amssymb, amsthm}
\usepackage{mathtools}

% --- Page layout ---
\usepackage[margin=1in]{geometry}

% --- Better lists ---
\usepackage{enumitem}
% --- Hyperlinks for references ---
\usepackage[colorlinks=true, linkcolor=blue, urlcolor=blue]{hyperref}

% Header: fancy layout like the rovided image
\usepackage{fancyhdr}
\setlength{\headheight}{26pt} % increase if fancyhdr warns
\pagestyle{fancy}
\fancyhf{} % clear header/footer
\rhead{\begin{tabular}{r} Mt.: 3052737\end{tabular}}
\chead{\begin{tabular}{c}Lösungen EA 2\end{tabular}}
\lhead{\begin{tabular}{l}61211 Analysis SS26 \end{tabular}}
\renewcommand{\headrulewidth}{0.8pt}
\fancyfoot[C]{- \thepage\ -}
% Ensure the plain page style (used by \maketitle) uses the same header

% --- Line spacing ---
\usepackage{setspace} % provides \doublespacing and \onehalfspacing

% --- Optional solutions environment ---
\newenvironment{solution}{
  \par\noindent\textbf{Solution:}\quad
}{\hfill$\blacksquare$\par\vspace{1em}}

% % --- Title info ---
% \title{Aufgabe \#1}
% \author{Your Name}
% \date{\today}
% 
\begin{document}
\raggedright
% Enable double (2.0) line spacing for the document body
\doublespacing
\noindent\textbf{Aufgabe 1} \\
\noindent\textbf{(a)} $f(x_1, x_2) = \frac{x_1^2 x_2}{1 + x_2^2}$. \\ Der Zähler und der Nenner sind Polynomfuktionen und daher stetig. Der Quotient ist stetig, da der Nenner ungleich null ist. Die Funktion ist daher auf ganz $\mathbb{R}^2$ stetig. \\
\bigskip
\noindent\textbf{(b)} \\ 
$f(x_1, x_2) = \begin{cases} \frac{x_1^2 + x_2^2}{x_1 x_2}, & x_1 \neq 0 \text{ und } x_2 \neq 0 \\ 0, & \text{sonst} \end{cases}$ \\
Für alle Punkte mit $x_1 \neq 0$ und $x_2 \neq 0$ ist $f$ als rationale Funktion auf ihrem Definitionsbereich stetig.\\
 Anwenden des Folgenkriteriums:  betrachte die Folge $(x_k) = (\frac{1}{k}, \frac{1}{k})$. Für $k \to \infty$ konvergiert $(x_k)$ gegen den Ursprung $(0,0)$. Die Funktion konvergiert gegen
 $$f \left(\frac{1}{k}, \frac{1}{k}\right) = \frac{(\frac{1}{k})^2 + (\frac{1}{k})^2}{(\frac{1}{k})(\frac{1}{k})} = \frac{\frac{2}{k^2}}{\frac{1}{k^2}} = 2.$$
Das das Folgenkriterium nicht erfüllt ist, ist die Funktion nicht stetig. \\
\bigskip
\noindent\textbf{(c)} \\
$f(x_1, x_2) = \begin{cases} \frac{x_1^2 + x_1 x_2 + x_2^2}{\sqrt{x_1^2 + x_2^2}}, & (x_1, x_2) \neq (0, 0) \\ 0, & \text{sonst} \end{cases}$ \\
Außerhalb des Ursprungs $(x_1, x_2) \neq (0,0)$: $f$ ist als Quotient und Verkettung stetiger Funktionen (Polynome und Quadratwurzel) stetig, da der Nenner ungleich null ist. \\
\renewcommand{\thefootnote}{\fnsymbol{footnote}} % *, †, ‡, …
Zur Untersuchung der Stetigkeit von \(f\) in \((0,0)\) betrachten wir für \((x_1,x_2)\neq(0,0)\) den Betrag
\begin{equation*}
|f(x_1, x_2)|
= \frac{|x_1^2 + x_1x_2 + x_2^2|}{\sqrt{x_1^2 + x_2^2}}
\le \frac{x_1^2 + |x_1x_2| + x_2^2}{\sqrt{x_1^2 + x_2^2}}
\stackrel{\footnotemark}{\le}
\frac{x_1^2 + \frac{x_1^2 + x_2^2}{2} + x_2^2}{\sqrt{x_1^2 + x_2^2}}
= \frac{3}{2}\sqrt{x_1^2 + x_2^2}.
\end{equation*}
Da
\[
\frac{3}{2}\sqrt{x_1^2 + x_2^2} \xrightarrow[(x_1,x_2)\to(0,0)]{} 0,
\]
folgt mit dem Sandwich-Theorem
\[
|f(x_1,x_2)| \xrightarrow[(x_1,x_2)\to(0,0)]{} 0.
\]
Somit gilt
\[
f(x_1,x_2) \xrightarrow[(x_1,x_2)\to(0,0)]{} 0 = f(0,0),
\]
also ist \(f\) in \((0,0)\) stetig.

\footnotetext{Hier wurde die Ungleichung \(|x_1x_2| \le \frac{x_1^2+x_2^2}{2}\) verwendet; vgl.\ den Beweis der Cauchy-Schwarz-Ungleichung (1.1.8).}
\end{document}